\typeout{NT FILE methodology.tex}%


\chapter{Methodology}
\label{cha:methodology}

In this chapter, we explain the methodology used in the development of this work. We begin 
by detailing how the data was obtained and the rationale behind its selection. Additionally, we 
present the data augmentation techniques applied to the dataset and describe the methodology 
employed in training the models.

\section{Data Collection}

The \gls{dgt} conducted an extensive aerial photogrammetric survey of Continental Portugal, 
resulting in a high-resolution orthophoto dataset covering the years 2018 to 2021. 
These orthophotos provide a valuable resource for geographic and environmental analysis, 
including housing fire risk assessment.

\subsection{Technical Specifications}
The orthophotos used in this study have the following characteristics:

\begin{itemize}
    \item \textbf{Spatial Resolution:} 25 cm/pixel.
    \item \textbf{Spectral Resolution:} Four-band imagery, including Red, Green, Blue (RGB), 
          and Near Infrared (NIR), enabling vegetation analysis.
    \item \textbf{Geographic Coverage:} Complete coverage of Continental Portugal.
    \item \textbf{Reference System:} 
        \begin{itemize}
            \item \textbf{EPSG:3763} - A projected coordinate system optimized for accurate 
                  distance measurements, ideal for mapping, surveying, and GIS applications in Portugal.
            \item \textbf{EPSG:4258} - A geographic coordinate system primarily used for geodetic 
                  applications and alignment across Europe.
        \end{itemize}
\end{itemize}

The orthophotos are publicly available through a Web Map Tile Service (WMTS), allowing for 
efficient data retrieval and integration into research applications. Our dataset consists of 
RGB and IRG orthophotos extracted from this open data source. To automate image retrieval, 
a Python module was developed to extract images using bounding boxes centered on 
the residential locations in the dataset. In total, 1930 RGB and IRG images were acquired.

For housing fire risk assessment, a 62.5-meter radius was defined to ensure that the analysis 
captures the area surrounding each house within at least 50 meters, with the additional 
12.5 meters accounting for the house itself. This approximation accounts for variability in 
house shape and size. Each extracted image is represented as a 500 x 500 pixel tile.

\begin{figure}[ht]
    \centering
    \begin{minipage}[b]{0.45\textwidth}
        \centering
        \includegraphics[width=\textwidth]{../Figures/houseExampleRGB.png}
        \caption{Example of an aerial photograph in \gls{rgb} format.}
        \label{fig:rgb}
    \end{minipage}
    \hfill
    \begin{minipage}[b]{0.45\textwidth}
        \centering
        \includegraphics[width=\textwidth]{../Figures/houseExampleIRG.png}
        \caption{Example of an aerial photograph in \gls{irg} format.}
        \label{fig:irg}
    \end{minipage}
    \caption{Comparison of aerial photographs of the same location in \gls{rgb} and \gls{irg} formats.}
    \label{fig:rgb-irg-comparison}
\end{figure}

\subsection{Ground Truth Collection}

The housing fire risk classifications provided by the \gls{adai} study were stored in a `.kmz` 
file format, viewable through Google Earth. This format allows for the visualization of the 
exact locations of isolated houses and their assigned fire risk levels. As illustrated in 
Figure \ref{fig:gearth_vis}, the classifications are color-coded: blue for low risk, green 
for moderate, yellow for high, orange for very high, and red for extreme risk.

\begin{figure}[ht]
    \centering
    \includegraphics[width=0.6\textwidth]{../Figures/GoogleEarthVis.png}
    \caption{Visualization of classified houses in Braga using Google Earth.}
    \label{fig:gearth_vis}
\end{figure}

To utilize this data, a Python script was developed to parse the `.kmz` file and extract the 
coordinates and corresponding risk levels for each house. The extracted information was saved 
into a `.csv` file, which was then used by the data collection pipeline to retrieve aerial 
imagery from the \gls{dgt} orthophoto service.

The final dataset consists of 1930 isolated houses across Portugal, categorized into five fire 
risk levels, as shown in Table~\ref{fig:data-distribution}.

\begin{table}[ht]
    \centering
    \renewcommand{\arraystretch}{1.2}
    \setlength{\tabcolsep}{12pt}
    \begin{tabular}{lc}
        \toprule
        \textbf{Fire Risk} & \textbf{Number of Houses} \\ 
        \midrule
        Low                &  25   \\ 
        Moderate           &  210  \\ 
        High               &  781  \\ 
        Very High          &  763  \\ 
        Extreme            &  151  \\ 
        \midrule
        \textbf{Total}     & \textbf{1930} \\ 
        \bottomrule
    \end{tabular}
    \caption{Distribution of housing fire risk classifications in the dataset.}
    \label{fig:data-distribution}
\end{table}
\subsection{Dataset Challenges and Simplification Strategy}

With the aerial imagery and fire risk labels collected, we next analyze the 
limitations and challenges inherent to this dataset.
These include interclass similarity, lighting variability, and 
class imbalance. To address them, a label simplification strategy 
was adopted.

\subsubsection*{Interclass Similarity and Visual Ambiguity}

Distinguishing between adjacent fire risk levels, such as High and Very High, is 
very difficult due to overlapping visual features like vegetation density and 
proximity to structures. These subtle distinctions are often hard to capture 
from a top-down perspective, as illustrated in Figure \ref{fig:veryhigh-vs-high}.

\begin{figure}[ht]
    \centering
    \begin{minipage}[b]{0.45\textwidth}
        \centering
        \includegraphics[width=\textwidth]{../Figures/VC83.png}
        \caption{Very High Risk House}
        \label{fig:vc83}
    \end{minipage}
    \hfill
    \begin{minipage}[b]{0.45\textwidth}
        \centering
        \includegraphics[width=\textwidth]{../Figures/VR95.png}
        \caption{High Risk House}
        \label{fig:vc102}
    \end{minipage}
    \caption{Very High and High Fire Risk Areas, respectively, in \gls{irg} format.}
    \label{fig:veryhigh-vs-high}
\end{figure}

\subsubsection*{Lighting Variability}

The aerial imagery in the dataset was collected under a range of environmental 
conditions, leading to inconsistencies in lighting, shadow, and contrast. 
Images captured at different times of day or under varying cloud cover can 
obscure important vegetation details. Figure \ref{fig:dark-bright-comparison} 
demonstrates this variation between images from Coimbra and Portalegre.

\begin{figure}[ht]
    \centering
    \begin{minipage}[b]{0.45\textwidth}
        \centering
        \includegraphics[width=\textwidth]{../Figures/darker_house_image.png}
        \caption{Example in lower brightness conditions.}
        \label{fig:darker}
    \end{minipage}
    \hfill
    \begin{minipage}[b]{0.45\textwidth}
        \centering
        \includegraphics[width=\textwidth]{../Figures/brighter_house_image.png}
        \caption{Example in higher brightness conditions.}
        \label{fig:brighter}
    \end{minipage}
    \caption{Aerial photographs under different lighting conditions.}
    \label{fig:dark-bright-comparison}
\end{figure}

\subsubsection*{Class Imbalance}

As shown in Table \ref{fig:data-distribution}, the dataset is heavily skewed, 
with the majority of samples labeled as High or Very High risk. This imbalance 
can cause deep learning models to become biased toward these dominant classes, 
resulting in poor performance on underrepresented labels like Low and Moderate.

\subsubsection*{Label Consolidation}

To mitigate the effects of both interclass ambiguity and imbalance, the original 
five fire risk levels were consolidated into two broader categories: 
\textit{Vegetation Managed} (Low, Moderate, High) and \textit{Vegetation Not Managed} 
(Very High, Extreme). This restructuring yielded a more balanced dataset, 1016 and 
914 samples per class, respectively, while preserving meaningful distinctions 
relevant to fire risk. Although this binary classification reduces granularity, 
it provides a more stable foundation for model training and evaluation. A return 
to multi-class classification may be viable in future work with a more evenly 
distributed dataset.


\section{Data Augmentation}

Given the dataset limitations, several data augmentation techniques were applied to 
increase training diversity and improve model generalization. These augmentations help 
the model learn robust features while mitigating overfitting—an important consideration 
given the relatively small dataset size and the complexity of the classification task.

\begin{table}[ht]
    \centering
    \renewcommand{\arraystretch}{1.2}
    \setlength{\tabcolsep}{12pt}
    \begin{tabular}{lc}
        \toprule
        \textbf{Augmentation} & \textbf{Parameters} \\ 
        \midrule
        Random Horizontal Flip & Probability: 50\% \\ 
        Random Vertical Flip & Probability: 50\% \\ 
        Random Rotation  & Range: $[-30^{\circ},30^{\circ}]$ \\ 
        Random Brightness Adjustment & Range: $[-20\%,20\%]$ \\ 
        \bottomrule
    \end{tabular}
    \caption{Data augmentation techniques applied to the dataset.}
    \label{tab:augmentations}
\end{table}

As shown in Table \ref{tab:augmentations}, four augmentation techniques were 
employed \cite{torchvision2016}. Horizontal and vertical flips were applied 
with a probability of 50\%, ensuring that vegetation and terrain appear in 
varied orientations relative to houses. Random rotations within a $[-30^{\circ}, 30^{\circ}]$ 
range further increased variability, allowing the model to recognize patterns from 
different perspectives. Brightness adjustments within a $[-20\%, 20\%]$ range 
were used to simulate lighting variations caused by different capture times or weather conditions.

The same augmentation techniques were applied to both RGB and IRG images; however, 
they were applied randomly during training, ensuring that no two training instances 
were identical. Some augmentations, particularly brightness adjustments, may behave 
differently between the two formats due to their distinct spectral properties, 
further contributing to dataset variability.

\subsection{Image Preprocessing}

Prior to training, all input images were normalized using modality-specific 
mean and standard deviation values. This preprocessing step ensures consistent 
scaling across inputs, which is crucial for stable and effective deep learning model training. 
Normalization was applied independently for RGB, IRG, and NDVI representations.

These statistics were computed using a custom Python script that iterated over 
the training dataset to calculate the mean and standard deviation across all pixels and channels. 
For RGB and IRG images, values were computed per channel; for NDVI, a single-channel vegetation index, 
a global mean and standard deviation were calculated. The resulting values were:

\begin{itemize}
    \item \textbf{RGB:} Mean = [0.4109, 0.3848, 0.3479], Std = [0.1963, 0.1623, 0.1376]
    \item \textbf{IRG:} Mean = [0.5130, 0.4105, 0.3839], Std = [0.1694, 0.1970, 0.1628]
    \item \textbf{NDVI:} Mean = 0.1542, Std = 0.1946
\end{itemize}

Normalization was applied after data augmentation and before feeding images into the model 
during training. For RGB+NDVI inputs, the NDVI channel was concatenated with the RGB tensor 
and normalized using its respective statistics. This modality-specific normalization helps 
account for differences in spectral characteristics and dynamic ranges across the various input types.


\section{Environment Setup}

To ensure a stable and reproducible training environment, we used Poetry \cite{poetry2023} for 
dependency management and a Docker \cite{docker2023} container with CUDA support for execution 
on the NOVA DI-Cluster. This setup enables efficient resource allocation while leveraging GPU 
acceleration for faster training.

Training was conducted on the \textit{alakazam} nodes of the cluster, equipped with the 
following hardware specifications:

\begin{itemize}
    \item \textbf{CPU:} Dual AMD EPYC 9124 (64 threads total)
    \item \textbf{GPU:} NVIDIA L4 (23GB VRAM, CUDA 12.0)
    \item \textbf{Memory:} 251GB RAM
\end{itemize}

The Docker container provides several command-line arguments that allow flexible control 
over the training process. The most frequently used parameters include:

\begin{itemize}
    \item \texttt{--train}, \texttt{--test}: Specify whether to train a new model or evaluate an existing one.
    \item \texttt{--algorithm}: Select the model architecture for training or evaluation.
    \item \texttt{--num\_epochs}, \texttt{--batch\_size}: Define the number of training epochs and batch size per iteration.
    \item \texttt{--lr}, \texttt{--wd}: Control the initial learning rate and weight decay for optimization.
    \item \texttt{--load\_weights}: Load pre-trained model weights for fine-tuning.
    \item \texttt{--save\_as}: Specify the output filename for storing the trained model and metrics.
    \item \texttt{--num\_classes}: Define the number of target classes for classification.
    \item \texttt{--ndvi}: Enable the use of a fourth channel for the NDVI index in input images.
\end{itemize}

By defining these parameters at runtime, the training pipeline maintains flexibility while 
ensuring consistency across different experiments. This setup facilitates efficient 
experimentation without requiring modifications to the core training script.


\section{Model Training and Evaluation}


This section outlines the methodology used for training, validation, and testing of the 
chosen models. The discussion of the results is deferred to a later chapter.

\subsection{Training Methodology}

The training process follows a transfer learning approach, in which pre-trained 
models are fine-tuned for the housing fire risk classification task. All models 
are initialized with ImageNet weights and optimized using the AdamW optimizer. 
The loss function used is Binary Cross-Entropy with Logits Loss (BCEWithLogitsLoss), 
which is well-suited for binary classification problems.

To prevent overfitting, an early stopping mechanism is applied: training halts 
if the validation loss does not improve for five consecutive epochs. Based on 
observed validation dynamics, training typically converges between 15 and 21 
epochs. The best-performing model, determined by lowest validation loss, is saved 
for final evaluation.

Due to GPU memory constraints, a batch size of 16 is used for ResNet-50, while 
EfficientNet-B5 and DenseNet-161 are trained with a batch size of 8.

To ensure robustness and reproducibility, all training, validation, and testing 
procedures are repeated three times. Results are averaged to provide a stable
and reliable performance estimate. Each run uses a fixed random seed, and the 
same seed is applied across all models within a given run (e.g., run 1 uses 
seed A for every model, run 2 uses seed B, and so on). This ensures consistency 
across architectures while capturing variability between runs.

Cross-validation was not employed in this work due to the high computational 
cost associated with training deep convolutional models, particularly in ensemble 
settings. While k-fold cross-validation can yield more robust performance estimates, 
it is computationally intensive when applied to deep learning pipelines. Instead, 
repeating experiments with fixed seeds and averaging the results provides a 
practical compromise, capturing performance variability while maintaining feasibility 
given available resources.


\subsection{Model Evaluation}

To assess model performance, we primarily rely on training and validation metrics.
Model selection is based on validation performance, specifically, the epoch 
with the lowest validation loss. Test accuracy and F1-score are reported 
solely to assess generalization ability and are not used during model 
selection or tuning.

The key evaluation metrics include:

\begin{itemize}
    \item \textbf{Validation Accuracy and Loss}: Tracked across epochs to monitor learning dynamics and select the best-performing checkpoint.
    \item \textbf{Training Accuracy and Loss}: Used to analyze the learning curve and detect signs of overfitting or underfitting.
    \item \textbf{Test Accuracy}: Reflects how well the selected model generalizes to unseen data.
    \item \textbf{Test F1-score}: A balanced measure of precision and recall, particularly relevant given the dataset's class imbalance.
\end{itemize}

All metrics, including training/validation loss and accuracy, 
are stored in JSON format for each epoch to enable further analysis 
and reproducibility.


Each model is trained using three different input configurations: RGB, 
IRG, and RGB+NDVI. Performance across these representations is compared 
to determine the most effective feature set for housing fire risk classification.


\subsection{Training Strategy and Hyperparameter Tuning}

The training process is structured into different stages to optimize model performance:

\begin{itemize}
    \item \textbf{Learning Rate Tuning:} Models are first trained with varying learning rates to determine the optimal initial rate.
    \item \textbf{Weight Decay Optimization:} Once an appropriate learning rate is established, different weight decay values are tested to improve generalization on that learning rate.
    \item \textbf{Gradual Unfreezing (DenseNet-161):} Due to its initial suboptimal performance, a gradual unfreezing strategy was employed for DenseNet-161, as explained in the next paragraph.
\end{itemize}

\subsubsection*{Gradual Unfreezing Strategy for DenseNet-161}

Fine-tuning pre-trained models can lead to instability and overfitting, particularly when working 
with smaller datasets. In initial training runs, unfreezing all layers of DenseNet-161 simultaneously 
risked updating deep feature representations too aggressively, potentially leading to feature 
degradation rather than effective adaptation. To mitigate this, we implemented a gradual unfreezing 
strategy, where layers were progressively unfrozen throughout training.

Instead of unfreezing all layers at once, which can lead to excessive parameter updates, layers were 
incrementally unfrozen based on the current epoch. The process follows these steps:

\begin{enumerate}
    \item \textbf{Epoch 0--2}: Only the classifier (fully connected layer) is trainable, while all convolutional layers remain frozen.
    \item \textbf{Epoch 3}: The shallowest layers (initial convolutional blocks) are unfrozen and fine-tuned.
    \item \textbf{Epoch 6}: Mid-level layers are unfrozen, allowing deeper feature representations to be fine-tuned.
    \item \textbf{Epoch 9}: The final layers are unfrozen, enabling full fine-tuning of the network.
\end{enumerate}

This approach is designed to allow the model to first adapt high-level features (classifier), then 
low-level spatial features, followed by mid-level and deep feature representations. By gradually 
unfreezing layers, we aim to retain valuable pre-trained features while progressively adapting the 
model to the housing fire risk classification task.

The impact of this strategy is further analyzed in the results section, where we evaluate its effect 
on model performance.



\section{Custom Classification Layer}

To adapt the pre-trained architectures for the housing fire risk classification task, the default 
fully connected (FC) layers of DenseNet-161, ResNet-50, and EfficientNet-B5 were 
replaced with a custom classification head. This modification enhances the network's 
ability to generalize while improving flexibility in learning complex patterns.

Table \ref{tab:classification_layer} presents the architecture 
of the custom classification layer.

\begin{table}[ht]
    \centering
    \renewcommand{\arraystretch}{1.3} % Adjust row height
    \setlength{\tabcolsep}{1pt} % Adjust column spacing
    \begin{tabular}{lcc}
        \toprule
        \textbf{Layer} & \textbf{Size} & \textbf{Activation} \\ 
        \midrule
        Fully Connected (FC1) & Input: extracted features, Output: 4096 & ReLU \\
        Dropout & Probability: 0.4 & - \\
        Fully Connected (FC2) & Input: 4096, Output: 4096 & ReLU \\
        Dropout & Probability: 0.4 & - \\
        Output Layer (Binary Classification) & Input: 4096, Output: 1 & Sigmoid \\
        \bottomrule
    \end{tabular}
    \caption{Structure of the custom classification layer.}
    \label{tab:classification_layer}
\end{table}

The custom classification layer consists of two fully connected (FC) layers with 
ReLU activation, followed by dropout regularization to improve generalization. 
The final layer outputs a single neuron with a Sigmoid activation function for 
binary classification.

The dropout rate of 0.4 was assumed based on standard deep learning practices, 
aiming to prevent overfitting while preserving meaningful feature representations. 

By integrating this custom classification layer, the model preserves the rich feature 
representations extracted by the deep feature extractor while enhancing adaptability 
to the housing fire risk classification task.


\section{Ensemble Model}

In search of improving classification performance, two ensemble models were explored: one using only 
\gls{rgb} input and another incorporating multiple input types. The motivation behind these 
ensembles was to leverage different deep learning architectures to enhance robustness 
and generalization.

\subsection{RGB-Only Ensemble}

The first ensemble model, illustrated in Figure \ref{fig:rgb_ensemble}, was designed 
using only \gls{rgb} images as input. While additional spectral information, such as near-infrared 
(NIR) bands or vegetation indices like \gls{ndvi}, can be beneficial, \gls{rgb} data is often the most 
accessible and widely available source of imagery. In many real-world scenarios, obtaining 
\gls{irg} images or computing \gls{ndvi} may not always be feasible due to sensor limitations or data 
availability. Therefore, this model serves as a baseline to evaluate how well fire risk 
can be classified using only standard \gls{rgb} imagery.

\begin{figure}[ht]
    \centering
    \includegraphics[width=0.8\textwidth]{../Figures/EnsembleRGBOnly.drawio.pdf}
    \caption{RGB-based ensemble model combining DenseNet-161, ResNet-50, and EfficientNet-B5.}
    \label{fig:rgb_ensemble}
\end{figure}

In this approach, \gls{rgb} images are processed through three different pre-trained architectures: 
DenseNet-161, ResNet-50, and EfficientNet-B5. The predictions from these models are then 
merged to produce the final classification result.


\subsection{Multi-Input Ensemble}


Although each individual model is trained on a single input type, the term \textit{Multi-Input Ensemble}
refers to the combination of models trained on diverse the inputs, \gls{rgb}, \gls{irg}, and \gls{rgb} + \gls{ndvi}. 
Rather than fusing these inputs at the model level, this approach makes use of their complementary spectral 
information during the ensemble stage.

Each input type is processed independently through three CNN architectures, DenseNet-161, ResNet-50, 
and EfficientNet-B5, resulting in a total of nine models. The ensemble aggregates their predictions 
to improve classification robustness across different vegetation conditions.


\begin{figure}[ht]
    \centering
    \includegraphics[width=0.7\textwidth]{../Figures/EnsembleNDVI.png}
    \caption{Multi-input ensemble model combining RGB, IRG, and RGB+NDVI inputs.}
    \label{fig:multi_input_ensemble}
\end{figure}


\subsection{Ensemble Merging Strategy and Evaluation}

To aggregate predictions from the different model branches within the ensemble, 
a majority voting strategy was employed. Each branch corresponds to a unique 
combination of input type (e.g., RGB, IRG, RGB+NDVI) and model architecture 
(DenseNet-161, ResNet-50, or EfficientNet-B5). Each of these branches 
independently produces a classification output, and the final ensemble 
decision is made based on the most frequently predicted class among them.

Majority voting was chosen over probability-based soft voting to reduce the 
influence of any single model's calibration or bias. This hard-voting method 
is particularly effective in binary classification settings where individual 
models have comparable performance, offering a simple yet robust decision 
fusion strategy.

By implementing these ensemble strategies, the goal was to assess whether 
incorporating additional spectral information (like NIR or NDVI) improves 
classification performance, and to determine how well fire risk can be predicted 
using only RGB imagery. This comparison is particularly relevant for real-world 
applications, where access to multi-spectral data may be limited.

\subsubsection{Ensemble Performance Metrics}

The ensemble models are evaluated using the following key metrics:

\begin{itemize}
    \item \textbf{Final Accuracy}: Measures the overall classification performance of the ensemble.
    \item \textbf{F1-score}: Ensures a balanced evaluation of precision and recall, accounting for class imbalances.
    \item \textbf{Confusion Matrix}: Provides insights into misclassification trends by comparing errors between independent models and the ensemble.
\end{itemize}

By combining the strengths of multiple models, the ensemble approach aims to improve classification 
consistency and mitigate individual model biases. 

\subsubsection{Analysis of High-Confidence Misclassifications}

To further evaluate the ensemble's reliability, we analyze cases where all models unanimously 
predict the same class, yet the ground truth label differs. These high-confidence misclassifications 
are particularly relevant, as they may reveal potential issues with the dataset, such as 
labeling inconsistencies or ambiguous features that make classification challenging. 

By investigating these cases, we aim to identify whether the errors stem from systematic model 
biases, limitations in the available features, or inaccuracies in the ground truth labels. 
Understanding these misclassifications helps refine both the dataset and the models, ensuring 
a more reliable and interpretable housing fire risk classification system.

\section{Methodology Overview Diagram}

To summarize the entire methodological process, Figure \ref{fig:methodology_pipeline} 
presents a visual overview of the pipeline developed in this study. It highlights the 
sequence from data acquisition and preprocessing through to model training, ensemble design, 
and final fire risk predictions.

\begin{figure}[ht]
    \centering
    \includegraphics[height=0.8\textheight]{../Figures/MethDiagram.pdf}
    \caption{End-to-end pipeline for housing fire risk classification using deep learning.}
    \label{fig:methodology_pipeline}
\end{figure}
